\begin{abstract}
Retrieving an object from a potentially hazardous area requires coordinated perception, mobility and manipulation. We present HazBot, a Raspberry Pi-centered mobile-manipulation architecture with four subsystems: LiDAR navigation, differential wheel drive, native camera inference and an articulated gripper. A direct CSI camera and separate drive, arm and sensor-controller paths support a seven-phase reference retrieval mission with explicit handoff prerequisites. Ball-based simulated control provides a nonhazardous task context; evaluation comprises three offline studies rather than physical mission trials. Exhaustive analysis of the isolated wheel-command branch covers 261,121 input pairs, identifies magnitude loss in direction-only translation and finds 1,020 mixed-zero pairs that request no new command. Two URDF links give a 0.200\,m shoulder-to-wrist bound; 68,921 joint samples characterize the ideal wrist envelope. The full reference coordinator produces no violations of the specified command contracts across 28,672 phase--observation pairs; removing individual guard groups exposes counterexamples. The contribution is an implementation-grounded architecture and coordination study that identifies gateway, feedback and calibration requirements for subsequent task-specific retrieval experiments.
\end{abstract}
